\documentclass[a4paper,10pt]{article}
\usepackage[utf8]{inputenc}
\usepackage[T1]{fontenc}
\usepackage[sc]{mathpazo}
\usepackage[scaled=0.88]{helvet}
\usepackage{amsmath,amssymb,mathrsfs,stmaryrd}
\usepackage{microtype}
\usepackage[a4paper,top=1.5cm,bottom=1.5cm,left=1.4cm,right=1.4cm,headheight=12pt,headsep=0.4cm,footskip=0.7cm]{geometry}
\usepackage{xcolor}
\definecolor{Ink}{HTML}{1C2533}\definecolor{Bleu}{HTML}{1F3A5F}\definecolor{BleuMoyen}{HTML}{4A6A8F}
\definecolor{BleuClair}{HTML}{9DB2C9}\definecolor{Fond}{HTML}{F2F4F7}\definecolor{Ocre}{HTML}{B07A26}
\definecolor{Brique}{HTML}{A23B2A}\definecolor{Gris}{HTML}{5E6672}\definecolor{Filet}{HTML}{C5CCD6}
\usepackage{multicol,enumitem,booktabs,array,needspace}
\usepackage[most]{tcolorbox}
\usepackage{tikz,pgfplots}
\pgfplotsset{compat=1.17,trig format plots=rad}
\usetikzlibrary{arrows.meta}
\usepackage{fancyhdr}
\usepackage[colorlinks=true,linkcolor=Bleu,bookmarksopen=true,bookmarksopenlevel=0]{hyperref}

\setlength{\parindent}{0pt}\setlength{\parskip}{0.25em}
\setlength{\columnsep}{0.55cm}\setlength{\columnseprule}{0.3pt}\renewcommand{\columnseprulecolor}{\color{Filet}}
\setlist{itemsep=0.05em,topsep=0.15em,parsep=0pt,leftmargin=1.1em}
\setlist[itemize]{label=\textcolor{BleuMoyen}{\scriptsize\textbullet}}
\setlist[enumerate]{label=\textsf{\scriptsize\color{Bleu}\arabic*.},leftmargin=1.3em}
\linespread{1.02}

\newcommand{\K}{\mathbb{K}}\newcommand{\R}{\mathbb{R}}\newcommand{\C}{\mathbb{C}}\newcommand{\Z}{\mathbb{Z}}
\newcommand{\Mn}[1]{\mathcal{M}_{n}(#1)}\newcommand{\rg}{\operatorname{rg}}\newcommand{\vect}{\operatorname{Vect}}
\newcommand{\Mat}{\operatorname{Mat}}\newcommand{\Id}{\operatorname{Id}}\newcommand{\com}{\operatorname{com}}
\newcommand{\vv}[1]{\overrightarrow{#1}}\newcommand{\ve}[1]{\vec{#1}}
\newcommand{\dt}[1]{\left|\begin{matrix}#1\end{matrix}\right|}
\newcommand{\HP}{{\sffamily\bfseries\scriptsize\textcolor{Brique}{[HP]}}}
\newcommand{\livre}[1]{{\sffamily\scriptsize\color{Gris}[#1]}}

% ---------- titres ----------
\newcommand{\chapitre}[2]{\clearpage\phantomsection\pdfbookmark[0]{#1 --- #2}{chap#1}\markboth{#1 --- #2}{}%
  \noindent\colorbox{Bleu}{\parbox{\dimexpr\linewidth-2\fboxsep}{\color{white}\sffamily\bfseries\Large\strut #1\quad
  {\normalsize\mdseries #2}}}\par\vspace{0.4em}}
\newcommand{\notion}[1]{\par\needspace{9\baselineskip}\vspace{0.5em}\phantomsection\pdfbookmark[1]{#1}{\detokenize{#1}}\noindent{\sffamily\bfseries\color{Bleu}\large #1}\par\nobreak
  \noindent{\color{Ocre}\rule{\linewidth}{0.8pt}}\par\nobreak\vspace{0.1em}}
\newcommand{\sous}[1]{\par\needspace{6\baselineskip}\vspace{0.25em}\noindent{\sffamily\bfseries\small\color{BleuMoyen}#1}\par\nobreak}

% ---------- encadrés ----------
\tcbset{base/.style={enhanced,arc=0pt,boxrule=0pt,left=4pt,right=3pt,top=2pt,bottom=2pt,before skip=0.35em,after skip=0.35em,
  fontupper=\small}}
\newtcolorbox{DEF}[1][]{base,colback=white,borderline west={1.4pt}{0pt}{Bleu},
  before upper={{\sffamily\bfseries\scriptsize\color{Bleu}DÉFINITION\ifx\relax#1\relax\else\ --- #1\fi}\par}}
\newtcolorbox{PROP}[1][]{base,colback=Fond,borderline west={2pt}{0pt}{Bleu},
  before upper={{\sffamily\bfseries\scriptsize\color{Bleu}\ifx\relax#1\relax PROPRIÉTÉ\else #1\fi}\par}}
\newtcolorbox{FORM}[1][]{base,colback=white,colframe=BleuMoyen,boxrule=0.5pt,
  before upper={{\sffamily\bfseries\scriptsize\color{BleuMoyen}\ifx\relax#1\relax FORMULES\else #1\fi}\par}}
\newtcolorbox{METH}[1]{base,colback=white,borderline north={0.8pt}{0pt}{Bleu},borderline south={0.8pt}{0pt}{Bleu},
  before upper={{\sffamily\bfseries\scriptsize\colorbox{Bleu}{\color{white}MÉTHODE}\ \color{Bleu}#1}\par\smallskip}}
\newtcolorbox{PIEGE}{base,colback=Brique!5,borderline west={1.4pt}{0pt}{Brique},
  before upper={{\sffamily\bfseries\scriptsize\color{Brique}PIÈGES}\par}}
\newtcolorbox{ASTU}{base,colback=Ocre!7,borderline west={1.4pt}{0pt}{Ocre},
  before upper={{\sffamily\bfseries\scriptsize\color{Ocre}ASTUCES}\par}}
\newtcolorbox{HPBOX}{base,colback=white,borderline west={1.4pt}{0pt}{Brique},
  before upper={{\sffamily\bfseries\scriptsize\color{Brique}HORS PROGRAMME}\par}}

\pagestyle{fancy}\fancyhf{}
\fancyhead[L]{\scriptsize\sffamily\color{Gris}\leftmark}
\fancyhead[R]{\scriptsize\sffamily\scshape\color{Gris}Le Collé PT --- fiche Algèbre linéaire}
\fancyfoot[C]{\scriptsize\sffamily\thepage}
\renewcommand{\headrulewidth}{0.3pt}

\pgfplotsset{mini/.style={axis lines=middle,axis line style={Filet},xtick=\empty,ytick=\empty,width=3.2cm,height=3.2cm,
  xmin=-1.15,xmax=1.15,ymin=-1.15,ymax=1.15,title style={font=\scriptsize\sffamily,yshift=-0.15cm},samples=80}}


\newcommand{\Ker}{\operatorname{Ker}}\renewcommand{\Im}{\operatorname{Im}}\newcommand{\tr}{\operatorname{tr}}\newcommand{\Sp}{\operatorname{Sp}}
\newcommand{\GLn}[1]{\mathrm{GL}_n(#1)}
\hypersetup{pdftitle={Le Collé — fiche de référence : algèbre linéaire}}
\begin{document}
\thispagestyle{empty}
\noindent\colorbox{Bleu}{\parbox{\dimexpr\linewidth-2\fboxsep}{\vspace{0.4em}
{\color{white}\sffamily\scriptsize\textls[200]{LE COLLÉ\ \ ·\ \ MATHÉMATIQUES PT\ \ ·\ \ CHAPITRE 1}}\par\smallskip
{\color{white}\sffamily\bfseries\LARGE Fiche de référence --- Algèbre linéaire}\par\smallskip
{\color{white!85!Bleu}\sffamily rappels, sommes directes, projecteurs, hyperplans, trace, sous-espaces stables, résultats du TD}\vspace{0.4em}}}

\vspace{0.6em}
{\small Le livret sert à apprendre, la présentation à comprendre, cette fiche à travailler et réviser. Mêmes numéros de parties et
mêmes notations que la présentation et le livret ; \livre{th.~1.5} renvoie au livret, \livre{TD~12} à l'exercice~12 du TD complet.
$\K=\R$ ou $\C$, $f\in\mathcal L(E)$, $\dim E=n$, $A\in\Mn{\K}$.}

\vspace{0.4em}
{\small\sffamily\noindent\begin{tabular}{@{}p{0.6\linewidth}@{}}
1.1 Rappels de première année \dotfill\pageref{n11}\\
1.2 Produit et somme de sous-espaces ; projecteurs \dotfill\pageref{n12}\\
1.3 Hyperplans et systèmes linéaires \dotfill\pageref{n13}\\
1.4 Trace \dotfill\pageref{n14}\\
1.5 Sous-espaces stables \dotfill\pageref{n15}\\
1.6 Quelle méthode ? Pièges ; résultats du TD \dotfill\pageref{n16}\\
\end{tabular}}

\vspace{0.6em}
\pdfbookmark[0]{Chapitre 1 --- Algèbre linéaire}{chap1}\markboth{Chapitre 1 --- Algèbre linéaire}{}
\begin{multicols}{2}\raggedcolumns
% ======================================================================
\notion{1.1 Rappels de première année}\label{n11}
\begin{FORM}[CALCUL MATRICIEL \livre{R.1--R.4}]
$(AB)_{i,j}=\sum_ka_{i,k}b_{k,j}$ ; $AB\neq BA$ en général ; $(AB)^T=B^TA^T$.\\
Binôme si $AB=BA$. $(AB)^{-1}=B^{-1}A^{-1}$ ; $AB=I_n\Rightarrow B=A^{-1}$.\\
Triangulaire inversible $\iff$ diagonale sans zéro. $\begin{pmatrix}a&b\\c&d\end{pmatrix}^{-1}=\frac1{ad-bc}\begin{pmatrix}d&-b\\-c&a\end{pmatrix}$.\\
$AX=B$ compatible : solutions $X_0+\Ker A$. Pivot : opérations sur les lignes.
\end{FORM}
\begin{FORM}[FAMILLES ET DIMENSION \livre{R.9--R.12}]
Base extraite, base incomplète ; $n+1$ vecteurs dans un espace engendré par $n$ : liés.\\
$\dim\K^n=n$, $\dim\K_n[X]=n+1$, $\dim\mathcal M_{n,p}=np$, $\dim\mathcal L(E,F)=\dim E\dim F$.\\
$\dim E=n$ : $n$ vecteurs, base $\iff$ libre $\iff$ génératrice.\\
$\rg(x_1,\dots,x_p)=\dim\vect(x_i)\leq p$ ($=p\iff$ libre). Degrés distincts : libre.\\
$F\subset E$ : $\dim F\leq\dim E$, égalité $\iff F=E$.
\end{FORM}
\begin{PROP}[RANG \livre{prop.~R.15}]
$\Ker u$, $\Im u$ sous-espaces ; $u$ injective $\iff\Ker u=\{0\}$.\\
$(e_j)$ génératrice $\Rightarrow\Im u=\vect\big(u(e_j)\big)$.\\
\emph{Théorème du rang} : $\dim\Ker u+\rg u=\dim E$.\\
$\rg(v\circ u)\leq\min(\rg u,\rg v)$ ; inchangé par composition avec un isomorphisme.\\
$\dim E=\dim F$ : injective $\iff$ surjective $\iff$ bijective.\\
$u$ déterminée par $u(e_i)$ ; injective / surjective / bijective $\iff\big(u(e_i)\big)$ libre / génératrice / base \livre{R.16}.\\
$u(x)=a$ : aucune solution, ou $x_0+\Ker u$ \livre{R.17}.
\end{PROP}
\begin{FORM}[MATRICES \livre{R.18--R.23}]
Colonne $j$ : coordonnées de $u(e_j)$ ; $Y=AX$ ; $\Mat(v\circ u)=\Mat(v)\Mat(u)$.\\
$\Im A=\vect(\text{colonnes})$ ; $\rg A^T=\rg A$ ; $\dim\Ker A=p-\rg A$.\\
$A\in\Mn{\K}$ inversible $\iff\Ker A=\{0\}\iff\rg A=n$.\\
$AX=B$ compatible $\iff B\in\Im A$ ; $A$ inversible : $X=A^{-1}B$.
\end{FORM}
\begin{FORM}[CHANGEMENT DE BASE \livre{prop.~R.21}]
$P$ = matrice de passage de $\mathcal B$ à $\mathcal B'$ (colonnes : $\mathcal B'$ dans $\mathcal B$), $X=PX'$ :
\[\Mat_{\mathcal B'}(f)=P^{-1}\Mat_{\mathcal B}(f)\,P.\]
$A$, $B$ \emph{semblables} : $B=P^{-1}AP$. Mêmes rang, trace (§1.4), déterminant (ch.~2), polynôme caractéristique (ch.~3).
\end{FORM}
\begin{METH}{Inverse, rang, noyau \livre{méth.~R.1, R.3}}
Inverse : résoudre $AX=Y$, ou pivot sur $(A\,|\,I_n)$.\\
Rang : échelonner, compter les pivots. Noyau : $AX=0$, contrôler $\dim=p-\rg A$. Image : colonnes des pivots.
\end{METH}
% ======================================================================
\notion{1.2 Produit et somme de sous-espaces ; projecteurs}\label{n12}
\begin{PROP}[PRODUIT, SOMME DIRECTE \livre{prop.~1.1, 1.4, th.~1.5}]
$\dim(E_1\times\dots\times E_p)=\sum\dim E_i$.\\
$F_1+\dots+F_p$ directe $\iff$ ($x_1+\dots+x_p=0$, $x_i\in F_i\Rightarrow$ tous nuls). $p=2$ : $F\cap G=\{0\}$.\\
Base adaptée : concaténation des bases ; $\dim\bigoplus F_i=\sum\dim F_i$.\\
$E=F\oplus G\iff$ deux parmi : $F\cap G=\{0\}$, $F+G=E$, $\dim F+\dim G=n$.\\
Grassmann : $\dim(F+G)=\dim F+\dim G-\dim(F\cap G)$.
\end{PROP}
\begin{PIEGE}
$p\geq3$ : intersections deux à deux nulles $\not\Rightarrow$ directe ($\vect(1,0)$, $\vect(0,1)$, $\vect(1,1)$).
\end{PIEGE}
\begin{FORM}[PROJECTEURS, SYMÉTRIES \livre{prop.~1.6}]
$p\circ p=p$ : projection sur $\Im p$ parallèlement à $\Ker p$, $E=\Im p\oplus\Ker p$, $\Mat=\mathrm{diag}(I_r,0)$.\\
$s\circ s=\Id$ : $E=\Ker(s-\Id)\oplus\Ker(s+\Id)$, $\Mat=\mathrm{diag}(I_r,-I_{n-r})$ ; $s=2p-\Id$.\\
Calcul de $p(x)$ sur $F$ parallèlement à $G=\vect(w)$ : $t$ tel que $x-tw\in F$, $p(x)=x-tw$ \livre{TD~10}.
\end{FORM}
% ======================================================================
\notion{1.3 Hyperplans et systèmes linéaires}\label{n13}
\begin{PROP}[HYPERPLANS \livre{prop.~1.8}]
Hyperplan : supplémentaire d'une droite ; dimension $n-1$ ; $=\Ker\varphi$, $\varphi$ forme linéaire non nulle ; équation
$a_1x_1+\dots+a_nx_n=0$, unique à un facteur près.\\
$u\notin H\Rightarrow E=H\oplus\vect(u)$.\\
$p$ hyperplans : intersection de dimension $\geq n-p$ ; sous-espace de dimension $n-p$ = $p$ équations.
\end{PROP}
\begin{METH}{Base $\leftrightarrow$ équations \livre{méth.~1.1}}
Équations $\to$ base : résoudre. Base $(v_1,\dots,v_k)$ $\to$ équations : conditions de compatibilité de $x=\sum\lambda_iv_i$ ;
il y en a $n-k$.\\
Système homogène de rang $r$ : solutions de dimension $n-r$ ; avec second membre : vide ou $X_0+\Ker A$.
\end{METH}
% ======================================================================
\notion{1.4 Trace}\label{n14}
\begin{FORM}[PROPRIÉTÉS \livre{prop.~1.11}]
$\tr A=\sum a_{i,i}$ : forme linéaire ; $\Ker\tr$ hyperplan de $\Mn{\K}$.\\
$\tr(AB)=\tr(BA)$ ; $\tr(P^{-1}AP)=\tr A$ ; $\tr f$ indépendante de la base.\\
Projecteur : $\tr p=\rg p$. Symétrie : $\tr s=\dim\Ker(s-\Id)-\dim\Ker(s+\Id)$.\\
$AB-BA=I_n$ impossible \livre{TD~11} (en dimension finie seulement, \livre{TD~21}).
\end{FORM}
\begin{PIEGE}
$\tr(AB)\neq\tr A\,\tr B$ ; $\tr(ABC)=\tr(CAB)$ mais $\neq\tr(BAC)$ en général.
\end{PIEGE}
% ======================================================================
\notion{1.5 Sous-espaces stables}\label{n15}
\begin{DEF}
$F$ stable : $f(F)\subset F$ ; induit $f_F\in\mathcal L(F)$. Vérifier $f(e_j)\in F$ sur une base de $F$.\\
$f\circ g=g\circ f\Rightarrow\Ker g$, $\Im g$ stables par $f$ ; $\Ker f$, $\Im f$ stables par $f$.
\end{DEF}
\begin{PROP}[BLOCS \livre{th.~1.15}]
Base adaptée à $F$ ($\dim F=p$) : $F$ stable $\iff\Mat(f)=\begin{pmatrix}A&B\\0&C\end{pmatrix}$, $A=\Mat(f_F)$.\\
$E=F\oplus G$, les deux stables $\iff\begin{pmatrix}A&0\\0&C\end{pmatrix}$ ; $\tr f=\tr A+\tr C$.\\
Droite $\vect(u)$ stable $\iff f(u)=\lambda u$ \livre{prop.~1.16} : point de départ de la réduction (ch.~3).
\end{PROP}
% ======================================================================
\notion{1.6 Quelle méthode ? Pièges ; résultats du TD}\label{n16}
{\small\renewcommand{\arraystretch}{1.12}
\begin{tabular}{@{}>{\raggedright\arraybackslash}p{0.34\linewidth}>{\raggedright\arraybackslash}p{0.6\linewidth}@{}}
\hline
\textbf{Question} & \textbf{Méthode}\\\hline
$\rg u$, $\dim\Ker u$ & base de l'image, théorème du rang\\
$E=F\oplus G$ ? & intersection nulle + dimensions\\
projecteur, symétrie ? & $f^2=f$, $f^2=\Id$ ; base adaptée\\
équations de $F$ & compatibilité de $x=\sum\lambda_iv_i$\\
trace & linéaire, $\tr(AB)=\tr(BA)$\\
$F$ stable ? & images d'une base de $F$ ; blocs\\\hline
\end{tabular}}
\begin{PIEGE}
$E=F\oplus G$ demande aussi la somme ou les dimensions. $\Ker f$ et $\Im f$ ne sont pas supplémentaires en général
($f^2=0$, \livre{TD~18}). Matrice de passage : colonnes = \emph{nouveaux} vecteurs dans l'\emph{ancienne} base.
\end{PIEGE}
{\small
\renewcommand{\arraystretch}{1.12}
\begin{tabular}{@{}>{\sffamily\scriptsize\bfseries\color{Bleu}}l>{\raggedright\arraybackslash}p{0.84\linewidth}@{}}
TD4 & $\{\sum x_i=0\}$ : hyperplan, base $(e_i-e_n)_{i<n}$.\\
TD6 & $\R^3=\{x+y+z=0\}\oplus\vect(1,1,1)$, $m=\frac{x+y+z}3$.\\
TD10 & projection sur $2x-y+z=0$ parallèlement à $\vect(1,0,1)$ : $t=\frac{2x-y+z}3$.\\
TD12 & $\tr(AE_{j,i})=a_{i,j}$ : $\tr(AM)=\tr(BM)\ \forall M\iff A=B$.\\
TD13 & rang de $\begin{pmatrix}a&b&1\\1&ab&1\\1&b&a\end{pmatrix}$ : $1$ si $a=1$ ; $2$ si $b=0$ ou $a=-2$ ; $3$ sinon.\\
TD14 & $A^2=A$, $\rg A=\tr A=2$ ; $\Im$ : $x-2y+z=0$.\\
TD17 & $\rg M\leq1\iff M=UV^T$ ; alors $M^2=\tr(M)M$.\\
TD19 & $M\mapsto AM-MA$ : $\Ker\supset\vect(I,A)$, $\Im\subset\Ker\tr$.\\
TD20 & $M^2-\tr(M)M+\det(M)I_2=0$ ; nilpotente $\iff\tr=\det=0$.\\
TD22 & $\tr f=0$ en dim.~2 : $f^2=\lambda\Id$ ; $f=\alpha s$ ou $\begin{pmatrix}0&1\\0&0\end{pmatrix}$.
\end{tabular}}
\end{multicols}
\end{document}
